\subsection{Evaluation}

\subsubsection{Advantages and Disadvantages}

\paragraph{Advantages}
\begin{itemize}
    \item \textbf{Strict Encapsulation:} Client code never accesses the internal node pointers or collection arrays directly~\cite{gof1994}.
    \item \textbf{Single Responsibility Principle:} Traversal logic and cursor bookkeeping are isolated inside the iterator class, keeping the container focused solely on data storage.
    \item \textbf{Open/Closed Principle:} New data structures or alternative traversal strategies (e.g., filtered, reverse, or breadth-first) can be added without modifying existing consumer loops.
    \item \textbf{Independent Concurrent Traversals:} Because each iterator maintains its own cursor state, multiple independent traversals can run across the same underlying container simultaneously.
    \item \textbf{Lazy Evaluation / Memory Efficiency:} Elements can be processed on-demand, enabling safe traversal of streaming data, infinite sequences, or large database cursors without exhausting memory.
\end{itemize}

\paragraph{Disadvantages}
\begin{itemize}
    \item \textbf{Over-Engineering for Simple Needs:} Implementing polymorphic interfaces, factory methods, and custom iterator classes introduces unnecessary complexity for simple, fixed-size structures.
    \item \textbf{Indirection Overhead:} Virtual function dispatch and dynamic memory allocations (\texttt{new} / \texttt{std::make\_unique}) introduce CPU and heap overhead compared to native index loops~\cite{stroustrup2013}.
\end{itemize}

\subsubsection{Iterator Invalidation and Object Lifetime Pitfalls}
In native C++, memory safety and object lifetimes are critical considerations when implementing custom iterators~\cite{stroustrup2013, meyers2014}:
\begin{itemize}
    \item \textbf{Dangling References / Lifetime Discrepancy:} A concrete iterator typically retains a raw reference (e.g., \texttt{const std::vector<Track>\&}) or a pointer to the aggregate container. If the aggregate object's lifetime ends before the iterator is destroyed (for example, if the aggregate was created as a temporary rvalue and discarded), subsequent calls to \texttt{hasNext()} or \texttt{next()} will dereference dangling memory, resulting in undefined behavior.
    \item \textbf{Structural Mutation during Traversal:} If the underlying container is structurally modified (e.g., inserting or appending elements via \texttt{push\_back()}) while an iterator is actively traversing it, \texttt{std::vector} may trigger internal heap reallocation~\cite{stroustrup2013}. This moves the underlying buffer to a new memory address, invalidating any active pointers or indices held by external iterators.
    \item \textbf{Architectural Mitigations:} To safeguard against invalidation in production systems, engineers can: (1) enforce strict scoped iteration where container lifetimes strictly enclose iterator lifetimes, (2) share ownership via \texttt{std::shared\_ptr<const Container>}, or (3) adopt snapshot-based iterators.
\end{itemize}

\subsubsection{High-Performance Optimization Techniques in C++}

\paragraph{Zero-Cost Abstractions and Stack-Allocated Iterators}
In performance-critical C++, virtual method calls and heap allocation can be bypassed entirely by utilizing stack-allocated value types that satisfy the C++ standard range idiom (\texttt{begin()} and \texttt{end()})~\cite{stroustrup2013, meyers2014}:

\begin{lstlisting}[language=C++, caption={Zero-Cost Stack-Allocated Iterator in C++}]
class HighPerformanceBuffer {
private:
    std::vector<Track> buffer;

public:
    // Lightweight Stack-Allocated Iterator
    struct Iterator {
        const Track* ptr;

        bool operator!=(const Iterator& other) const { 
            return ptr != other.ptr; 
        }
        const Track& operator*() const { return *ptr; }
        Iterator& operator++() { ++ptr; return *this; }
    };

    Iterator begin() const { return Iterator{buffer.data()}; }
    Iterator end() const { 
        return Iterator{buffer.data() + buffer.size()}; 
    }
};

// Enables compiler-inlined, zero-cost range-based for loops:
void renderAll(const HighPerformanceBuffer& buff) {
    for (const auto& track : buff) {
        std::cout << track.title << "\n";
    }
}
\end{lstlisting}

\paragraph{Internal Iterators via Higher-Order Functions}
Unlike external iterators where the client drives traversal manually via loops, internal iterators encapsulate loop execution inside the aggregate and apply a callback to each item:

\begin{lstlisting}[language=C++, caption={Internal Iterator via Functional Callback}]
class SongCollection {
    std::vector<Track> songs;
public:
    template <typename Callback>
    void forEach(Callback action) const {
        for (const auto& song : songs) {
            action(song);
        }
    }
};
\end{lstlisting}

\subsubsection{Concurrency, Mutation Handling, and C++ Standards}
Modifying a collection while an iterator is actively traversing it leads to race conditions, invalid iterators, or data races~\cite{williams2019}. Table~\ref{tab:thread_safe_iteration} contrasts standard synchronization techniques:

\begin{table}[htbp]
\centering
\caption{Comparison of Thread-Safe Iteration Strategies}
\label{tab:thread_safe_iteration}
\begin{tabular}{@{}p{3.6cm} p{4cm} p{3.8cm} p{3.2cm}@{}}
\toprule
\textbf{Technique} & \textbf{Mechanism} & \textbf{Advantages} & \textbf{Disadvantages} \\ \midrule
\textbf{Shared Mutex} (\texttt{std::shared\_mutex}) & Uses shared locks for reading/iteration; exclusive locks for writing~\cite{williams2019}. & Simple; minimal memory overhead. & High read load can starve pending write operations. \\ \midrule
\textbf{Defensive Snapshot} & Iterator duplicates the underlying collection at creation time. & Eliminates reader-writer conflicts completely. & Allocation and copying overhead proportional to collection size. \\ \midrule
\textbf{Atomic Pointer Swap (Copy-on-Write)} & Traversals read an immutable shared snapshot; writes clone and atomically swap pointers~\cite{williams2019, iso_cpp20}. & Non-blocking read traversal; reader threads never block on writer threads. & Lock-free capability is implementation-defined; write-amplification cost. \\ \bottomrule
\end{tabular}
\end{table}

\begin{lstlisting}[language=C++, caption={Copy-on-Write Snapshot Traversal in C++20}]
#include <iostream>
#include <vector>
#include <memory>
#include <atomic>

// Note: std::atomic<std::shared_ptr<T>> is standardized in C++20 (ISO/IEC 14882:2020).
// In C++11/C++14/C++17, the equivalent was std::atomic_load/store free functions.
class LockFreeTrackCatalog {
private:
    std::atomic<std::shared_ptr<const std::vector<Track>>> catalog_ptr;

public:
    LockFreeTrackCatalog() {
        catalog_ptr.store(std::make_shared<const std::vector<Track>>());
    }

    // Read traversal: readers acquire an immutable snapshot without locking mutexes.
    void traverse() const {
        std::shared_ptr<const std::vector<Track>> snapshot = catalog_ptr.load();
        for (const auto& track : *snapshot) {
            std::cout << track.title << "\n";
        }
    }

    // Write operation using atomic compare-and-swap (CAS)
    void addTrack(const Track& track) {
        std::shared_ptr<const std::vector<Track>> old_data;
        std::shared_ptr<const std::vector<Track>> new_data;

        do {
            old_data = catalog_ptr.load();
            auto copy = std::make_shared<std::vector<Track>>(*old_data);
            copy->push_back(track);
            new_data = copy;
        } while (!catalog_ptr.compare_exchange_weak(old_data, new_data));
    }

    // Utility to verify platform lock-free guarantee at runtime
    bool checkLockFree() const {
        return catalog_ptr.is_lock_free();
    }
};
\end{lstlisting}

\begin{quote}
\small \textbf{Important ISO C++ Standard Note (Lock-Freedom):} According to the C++ Working Draft~\cite{iso_cpp20, williams2019}, whether \texttt{std::atomic<std::shared\_ptr<T>>} is strictly lock-free on a given platform is \textbf{implementation-defined}. It relies on hardware support for double-word compare-and-swap (DWCAS / \texttt{CMPXCHG16B}) or small internal table-based locks. Developers can check \texttt{catalog\_ptr.is\_lock\_free()} at runtime or query \texttt{std::atomic<std::shared\_ptr<T>>::is\_always\_lock\_free} at compile time.
\end{quote}